\begin{table}[ht]
\centering
\caption{Distribuição dos ciclos de decisão por classe analítica.}
\label{tab:classes-ciclos}
\begin{tabular}{lr}
\toprule
\textbf{Classe} & \textbf{Quantidade} \\
\midrule
Ciclos estruturais e de generalização & 39 \\
Ciclos de consistência teórica e metodológica & 27 \\
Ciclos de refinamento visual e operacional & 47 \\
\bottomrule
\end{tabular}
\end{table}
\subsubsection{Ciclos estruturais e de generalização}
\begin{landscape}
\begingroup
\setlength{\tabcolsep}{3pt}
\small
\begin{longtable}{p{1.0cm}p{1.9cm}p{4.0cm}p{5.0cm}p{6.2cm}p{5.7cm}}
\caption{Ciclos estruturais e de generalização incorporados em versões consistentes.}\\
\toprule
\textbf{Ciclo} & \textbf{Data} & \textbf{Objetivo} & \textbf{Proposta inicial} & \textbf{Intervenção do pesquisador} & \textbf{Decisão consolidada} \\
\midrule
\endfirsthead
\toprule
\textbf{Ciclo} & \textbf{Data} & \textbf{Objetivo} & \textbf{Proposta inicial} & \textbf{Intervenção do pesquisador} & \textbf{Decisão consolidada} \\
\midrule
\endhead
C02 & 2026-06-30 & Ampliar o corpus de problemas & Carregamento de exemplos isolados & Agrupar por categoria e sortear situações curadas & Repositório organizado por categoria \\
\addlinespace[2pt]
C04 & 2026-07-01 & Oferecer retorno durante o arraste & Conjunto de efeitos visuais acoplados & Organizar proximidade, destaque, passagem, affordance e atração magnética & Strategy para estilos de proximidade \\
\addlinespace[2pt]
C08 & 2026-07-02 & Distinguir autoria das ações & Agente tratado de modo incompleto & Preservar S para sujeito e C para computador & Log com autoria explícita \\
\addlinespace[2pt]
C09 & 2026-07-02 & Distinguir ações sem finalidade matemática & Ações neutras e de interface agrupadas & Separar ações neutras de ações de interface & Taxonomia do log mais precisa \\
\addlinespace[2pt]
C11 & 2026-07-05 & Sincronizar diagramas em dois passos & Representações independentes & Usar o estado final do passo 1 como estado inicial do passo 2 & Sincronização entre representações encadeadas \\
\addlinespace[2pt]
C12 & 2026-07-05 & Organizar funcionalidades de proximidade & Recursos dispersos no código & Mover recursos para Scaffolding/Proximidade & Responsabilidades mais localizadas \\
\addlinespace[2pt]
C13 & 2026-07-13 & Oferecer quatro idiomas & Tradução com risco de alterar estado & Restringir troca de idioma à camada textual & Modelagem preservada entre idiomas \\
\addlinespace[2pt]
C15 & 2026-07-13 & Associar explicações à resolução & Tela genérica desvinculada da tentativa & Organizar tarefa matemática, interação, pesquisador e fotografia & Análise vinculada à tentativa \\
\addlinespace[2pt]
C16 & 2026-07-13 & Mostrar tentativas ao pesquisador & Exibir identificador técnico completo & Apresentar número ordinal por situação & Rastreabilidade legível sem perder unicidade \\
\addlinespace[2pt]
C17 & 2026-07-13 & Registrar formas simbólicas & Campo textual livre & Usar catálogo, códigos e construtor controlado & Invariantes normalizados \\
\addlinespace[2pt]
C18 & 2026-07-13 & Associar invariantes às ações & Registrar invariante apenas ao final & Repetir os quatro campos em cada ação & Linhas autossuficientes para análise \\
\addlinespace[2pt]
C21 & 2026-07-13 & Montar explicações por categoria & Listas fixas de campos & Gerar um bloco por elemento semântico da categoria & Renderização dinâmica \\
\addlinespace[2pt]
C23 & 2026-07-13 & Reduzir confusão entre campos & Formulário amplo com rolagem & Mostrar somente elementos válidos da categoria & Curadoria dinâmica sem campos irrelevantes \\
\addlinespace[2pt]
C24 & 2026-07-13 & Definir a fonte dos elementos semânticos & Fonte global indistinta & Usar a categoria apenas como fonte da tarefa matemática & Separação entre domínio matemático e interação \\
\addlinespace[2pt]
C25 & 2026-07-13 & Distinguir tipos e consequências das ações & Ação genérica & Separar MODELAGEM/NAVEGABILIDADE e natureza/efeito & Log analiticamente mais preciso \\
\addlinespace[2pt]
C26 & 2026-07-14 & Mostrar mesma solução em categorias distintas & Painéis independentes ou cinco colunas & Exigir tabela com quatro colunas e categoria sobre o enunciado & Novo modo comparativo tabular \\
\addlinespace[2pt]
C32 & 2026-07-14 & Reduzir a concentração de responsabilidades em Main.java & Dez classes gráficas mantidas como classes internas estáticas & Extrair os elementos para gerard.campoaditivo.diagrama.elementos, centralizar o tema dos menus e acrescentar regressão headless de geometria e desenho & Elementos gráficos passam a ter arquivos próprios com compilação e teste de fumaça automatizado \\
\addlinespace[2pt]
C34 & 2026-07-14 & Permitir uso por pessoa leiga sem configuração manual & Script dependente de Java previamente instalado & Empacotar distribuição com instalação automática, runtime compatível, ícone, atalhos e opção de jpackage & Distribuição reduz dependências técnicas do usuário final \\
\addlinespace[2pt]
C45 & 2026-07-15 & Vincular o controle à quantidade do termo desconhecido & Atualizar apenas o ponto e o cartão & Fazer a barra do referendo aumentar ou diminuir conforme o controle & Uma alteração no eixo propaga-se à representação da quantidade desconhecida \\
\addlinespace[2pt]
C52 & 2026-07-15 & Generalizar a fonte semântica para todas as categorias & Manter regras específicas dispersas por tela & Centralizar papéis, valores, sinais, incógnita e participantes na curadoria & Uma fonte semântica única alimenta todas as categorias aditivas \\
\addlinespace[2pt]
C60 & 2026-07-15 & Propagar alterações do eixo vertical & Atualizar apenas o gráfico de barras & Atualizar também Vergnaud e o eixo do número relativo & Eixo vertical torna-se fonte válida de alteração semântica \\
\addlinespace[2pt]
C61 & 2026-07-15 & Unificar o valor relativo entre representações & Manter atualizações parciais em uma direção & Sincronizar Vergnaud, barras e eixos em todas as ações e nos dois sentidos & Valor relativo passa a ter comportamento bidirecional consistente \\
\addlinespace[2pt]
C62 & 2026-07-15 & Permitir alterar o cartão do valor relativo & Tratar o cartão como exibição passiva & Propagar a edição para Vergnaud, barras e eixos & Edição textual participa do mesmo fluxo reativo \\
\addlinespace[2pt]
C64 & 2026-07-15 & Propagar manipulações das unidades & Mover quadradinhos sem alterar o modelo semântico & Atualizar Vergnaud após movimentação ou remoção de unidades & Venn e Vergnaud passam a compartilhar alterações \\
\addlinespace[2pt]
C66 & 2026-07-15 & Sincronizar a coleção Todo & Manter a resultante incompleta ou independente & Vincular a coleção Todo ao valor modelado e às manipulações bidirecionais & Todo torna-se manipulável e sincronizado \\
\addlinespace[2pt]
C67 & 2026-07-15 & Eliminar cadeias locais de sincronização & Atualizar representações diretamente umas pelas outras & Criar estado semântico único com snapshot, origem e proteção contra recursão & Todas as representações tornam-se projeções do mesmo estado \\
\addlinespace[2pt]
C73 & 2026-07-15 & Registrar problemas encontrados pelo usuário & Depender de descrição externa ao sistema & Adicionar formulário de relato com contexto e persistência local & Relatos passam a ser vinculados ao estado da atividade \\
\addlinespace[2pt]
C75 & 2026-07-15 & Incluir o estado representacional no relato & Registrar apenas texto e categoria & Adicionar representações exibidas e extensões futuras da curadoria & Relatos passam a descrever o contexto visual da falha \\
\addlinespace[2pt]
C85 & 2026-07-15 & Eliminar travamento ao salvar ou fechar & Executar gravações intermediárias duplicadas & Unificar o fluxo de salvar e fechar e tratar divergências numéricas & Fechamento passa a realizar uma única gravação segura \\
\addlinespace[2pt]
C90 & 2026-07-15 & Responder a posicionamentos semanticamente incompatíveis & Exibir somente uma anotação genérica & Combinar tremor leve, som sutil e tooltip contextual em todas as categorias & Erro semântico recebe feedback multissensorial que desaparece após correção \\
\addlinespace[2pt]
C91 & 2026-07-16 & Oferecer orientação situada nas representações & Concentrar a ajuda em instruções gerais & Adicionar interrogações no texto, Vergnaud e diagrama complementar com três intenções de ajuda & Usuário pode pedir dúvida, continuidade ou próximo passo na área relevante \\
\addlinespace[2pt]
C98 & 2026-07-16 & Permitir criação incremental de unidades nos agrupamentos & Manter quadradinhos apenas por reconstrução automática ou arraste existente & Adicionar controle vetorial junto a cada coleção, registrar o clique e sincronizar a nova quantidade & Coleções tornam-se extensíveis pelo usuário sem criar uma segunda fonte de verdade \\
\addlinespace[2pt]
C100 & 2026-07-16 & Permitir remoção incremental de unidades visíveis & Manter apenas o controle de adição ou depender da tecla Delete & Adicionar controle vetorial de subtração condicionado à existência de quadradinhos, preservar textos remanescentes e sincronizar a cardinalidade & Coleções passam a permitir aumento e redução explícitos sem criar fonte de verdade paralela \\
\addlinespace[2pt]
C105 & 2026-07-16 & Impedir que a manipulação direta ultrapasse o valor semântico curado & Permitir adição ilimitada de quadradinhos e apenas sincronizar a quantidade resultante & Usar o valor curado ou derivado como limite e aplicar tremor, som e tooltip ao atingi-lo & Coleções manipuláveis respeitam a quantidade da situação e não produzem cardinalidades semanticamente impossíveis \\
\addlinespace[2pt]
C107 & 2026-07-16 & Expressar criação e remoção de unidades como capacidades arquiteturais explícitas & Manter controles e tela acoplados diretamente a CirculoVenn e a verificações concretas & Definir interfaces-base e especializadas, centralizar estado comum em classe abstrata e operar implementações por polimorfismo & Representações declaram capacidades por contratos; controles deixam de depender da classe concreta e novas implementações podem variar sem condicionais por categoria \\
\addlinespace[2pt]
C112 & 2026-07-16 & Corrigir o incremento do Referendo na comparação & Usar o índice visual como índice do estado semântico & Introduzir contrato explícito entre ordem visual e papéis semânticos por categoria & Cada clique altera uma unidade no papel correto e recalcula o valor relativo \\
\addlinespace[2pt]
C116 & 2026-07-16 & Dar continuidade física ao deslocamento sem perder precisão & Copiar instantaneamente as coordenadas do cursor & Introduzir controlador de mola amortecida, origem fantasma e conclusão exata na soltura & O objeto persegue o cursor com atraso limitado e valida na coordenada real de soltura \\
\addlinespace[2pt]
C117 & 2026-07-16 & Manter os valores conhecidos do texto coerentes com os diagramas & Atualizar apenas Vergnaud e a representação complementar & Criar contrato comum para marcadores textuais e sincronizador pelo estado compartilhado & Texto manipulável acompanha o estado sem alterar a situação curada \\
\addlinespace[2pt]
C120 & 2026-07-16 & Unificar regras de adição e remoção de quantidades positivas & Controlar botões por eventos locais e estados compartilhados entre agrupamentos & Calcular cada controle pelo estado semântico, limites curados, piso zero e relação aditiva & Texto, Vergnaud, coleções e barras permanecem sincronizados sem alterar a curadoria \\
\addlinespace[2pt]
\bottomrule
\end{longtable}
\endgroup
\end{landscape}
\subsubsection{Ciclos de consistência teórica e metodológica}
\begin{landscape}
\begingroup
\setlength{\tabcolsep}{3pt}
\small
\begin{longtable}{p{1.0cm}p{1.9cm}p{4.0cm}p{5.0cm}p{6.2cm}p{5.7cm}}
\caption{Ciclos de consistência teórica e metodológica incorporados em versões consistentes.}\\
\toprule
\textbf{Ciclo} & \textbf{Data} & \textbf{Objetivo} & \textbf{Proposta inicial} & \textbf{Intervenção do pesquisador} & \textbf{Decisão consolidada} \\
\midrule
\endfirsthead
\toprule
\textbf{Ciclo} & \textbf{Data} & \textbf{Objetivo} & \textbf{Proposta inicial} & \textbf{Intervenção do pesquisador} & \textbf{Decisão consolidada} \\
\midrule
\endhead
C01 & 2026-06-30 & Preservar as representações formais & Inserção de elementos gráficos auxiliares & Retirar elementos tracejados e estranhos ao modelo & Regra permanente de fidelidade representacional \\
\addlinespace[2pt]
C03 & 2026-07-01 & Permitir manipulação direta do enunciado & Arraste amplo de elementos & Restringir aos números e ao ponto de interrogação & Modelagem dependente da ação explícita do usuário \\
\addlinespace[2pt]
C05 & 2026-07-01 & Representar sinais de transformação & Digitação de número negativo & Separar valor e sinal em controles distintos & Valor e sinal tornaram-se dimensões independentes \\
\addlinespace[2pt]
C07 & 2026-07-01 & Editar elementos do diagrama & Edição externa ao diagrama & Usar duplo clique sobre quadrados e círculos & Edição direta incorporada como estilo \\
\addlinespace[2pt]
C19 & 2026-07-13 & Registrar dificuldade percebida & Classificação binária ou do pesquisador & Usar Fácil/Intermediária/Difícil pelo usuário & Autorrelato com três níveis \\
\addlinespace[2pt]
C27 & 2026-07-14 & Separar categoria de instância & Diagrama da categoria preenchido e interativo & Torná-lo vazio, estático e não interativo & Distinção explícita entre estrutura e instância \\
\addlinespace[2pt]
C33 & 2026-07-14 & Construir a situação com blocos corretos e blocos de outras categorias & Usar alternativas de outras categorias sem controlar seu afastamento semântico & Exigir distância semântica controlada: manter contexto, personagens, objeto, valores, vocabulário e forma linguística aproximada, variando principalmente a relação entre quantidades & Nova aba exige interpretação do diagrama sem reduzir a tarefa a um jogo dos sete erros \\
\addlinespace[2pt]
C36 & 2026-07-15 & Substituir a representação inadequada de Venn na comparação de medidas & Manter a estrutura circular utilizada nas demais categorias & Adotar duas barras de quantidades com unidades manipuláveis e valor relativo explícito & Comparação passa a tornar visíveis referido, referendo e valor relativo \\
\addlinespace[2pt]
C37 & 2026-07-15 & Sincronizar o gráfico com os papéis curados & Usar a ordem superficial dos números do enunciado & Exigir referido, referendo e valor relativo como fonte semântica do gráfico & O gráfico deixa de trocar os papéis das quantidades \\
\addlinespace[2pt]
C39 & 2026-07-15 & Explicitar pessoas e objetos envolvidos no enunciado & Participantes dependiam de extração automática do texto & Adicionar personagem\_1, personagem\_2 e personagem\_3 como metadados livres e persistentes por versão linguística & Entidades do enunciado passam a ser registradas explicitamente sem depender apenas de inferência automática \\
\addlinespace[2pt]
C41 & 2026-07-15 & Tornar o valor relativo uma relação observável & Mostrar somente uma chave e um texto de diferença & Usar escala de zero a n, ponto de controle e nomes dos participantes sob as barras & A relação pode ser explorada sem perder os papéis semânticos \\
\addlinespace[2pt]
C47 & 2026-07-15 & Evidenciar a parte comum entre referido e referendo & Usar a mesma cor neutra para todas as unidades & Colorir em vermelho suave, de baixo para cima, a mesma quantidade de unidades nas duas barras & A quantidade comum torna-se visível e as unidades excedentes permanecem diferenciadas \\
\addlinespace[2pt]
C49 & 2026-07-15 & Evitar representação inadequada na transformação de medidas & Manter o diagrama de Venn ao lado de Vergnaud & Ocultar o Venn nessa categoria e ampliar a área do diagrama formal & Transformação de medidas deixa de exibir representação complementar inadequada \\
\addlinespace[2pt]
C50 & 2026-07-15 & Preservar a modelagem efetiva do usuário & Preencher inicialmente o gráfico com valores curados & Iniciar as barras vazias até o posicionamento explícito no diagrama de Vergnaud & A representação mostra apenas valores efetivamente modelados \\
\addlinespace[2pt]
C51 & 2026-07-15 & Corrigir correspondência entre Vergnaud e barras & Associar elementos pela ordem visual & Mapear referido, valor relativo e referendo pelos papéis curados & Rótulos e valores deixam de ser trocados na comparação \\
\addlinespace[2pt]
C53 & 2026-07-15 & Evitar antecipação de valores não modelados & Preencher simultaneamente as duas barras & Preencher cada barra somente quando o papel correspondente for posicionado & Cada quantidade visual corresponde a uma ação explícita do usuário \\
\addlinespace[2pt]
C65 & 2026-07-15 & Evitar representação antecipada & Carregar o Venn com os valores curados & Iniciar vazio e preencher somente após ações no diagrama formal & Venn passa a refletir apenas a modelagem construída \\
\addlinespace[2pt]
C68 & 2026-07-15 & Representar integralmente o valor do Todo & Limitar a coleção a doze unidades & Usar grade adaptativa sem truncar a quantidade semântica & Quantidade desenhada passa a coincidir com a quantidade matemática \\
\addlinespace[2pt]
C70 & 2026-07-15 & Atualizar personagens após a curadoria & Manter rótulos antigos até nova renderização & Reaplicar imediatamente os nomes nos nós já exibidos & Diagrama corrente passa a refletir a curadoria sem recarregamento \\
\addlinespace[2pt]
C86 & 2026-07-15 & Manter uma única fonte para dados semânticos & Permitir editar semântica em cada tradução & Bloquear dados semânticos nas traduções e herdá-los da original & Somente a versão original define a semântica da situação \\
\addlinespace[2pt]
C87 & 2026-07-15 & Eliminar informação redundante & Manter o campo subtipo editável & Retirar o campo e preservar compatibilidade interna & Curadoria deixa de duplicar estrutura já definida por outros campos \\
\addlinespace[2pt]
C88 & 2026-07-15 & Aplicar a regra ao trocar o idioma & Bloquear somente quando a linha já era tradução & Atualizar imediatamente o estado editável ao selecionar outro idioma & Semântica permanece protegida durante toda a edição da tradução \\
\addlinespace[2pt]
C94 & 2026-07-16 & Preservar as divisões da tela antes da escolha da categoria & Ocultar painéis juntamente com enunciado e diagramas & Distinguir estrutura visual estável de conteúdo educativo transitório e manter os três painéis vazios & A tela mantém sua organização espacial e carrega somente os conteúdos após a seleção da categoria \\
\addlinespace[2pt]
C96 & 2026-07-16 & Vincular a análise qualitativa a uma tentativa efetivamente iniciada & Permitir a abertura da análise logo após o carregamento da situação & Exigir ao menos um posicionamento atual no diagrama de Vergnaud e ocultar novamente após restauração & Análises e fotografias passam a possuir uma ação de modelagem como referente empírico \\
\addlinespace[2pt]
C118 & 2026-07-16 & Preservar a função semântica da interrogação & Substituir o ponto de interrogação pelo valor derivado & Sincronizar somente valores conhecidos e manter a incógnita original visível & O texto não revela automaticamente a resposta e mantém o papel desconhecido \\
\addlinespace[2pt]
C119 & 2026-07-16 & Impedir adição de unidades antes da modelagem em Vergnaud & Liberar o controle complementar assim que a situação fosse carregada & Condicionar a adição à explicitação dos elementos semânticos no diagrama formal & A representação complementar deixa de antecipar a modelagem formal \\
\addlinespace[2pt]
C122 & 2026-07-16 & Impedir que a escolha do sinal produza quantidades negativas & Propagar o sinal relativo e limitar visualmente a quantidade resultante & Validar a relação antes da propagação, permitir sinal apenas na transformação ou no valor relativo e preservar quantidades não negativas & Escolhas incompatíveis são bloqueadas sem alterar dados curados nem gerar estados negativos \\
\addlinespace[2pt]
\bottomrule
\end{longtable}
\endgroup
\end{landscape}
\subsubsection{Ciclos de refinamento visual e operacional}
\begin{landscape}
\begingroup
\setlength{\tabcolsep}{3pt}
\small
\begin{longtable}{p{1.0cm}p{1.9cm}p{4.0cm}p{5.0cm}p{6.2cm}p{5.7cm}}
\caption{Ciclos de refinamento visual e operacional incorporados em versões consistentes.}\\
\toprule
\textbf{Ciclo} & \textbf{Data} & \textbf{Objetivo} & \textbf{Proposta inicial} & \textbf{Intervenção do pesquisador} & \textbf{Decisão consolidada} \\
\midrule
\endfirsthead
\toprule
\textbf{Ciclo} & \textbf{Data} & \textbf{Objetivo} & \textbf{Proposta inicial} & \textbf{Intervenção do pesquisador} & \textbf{Decisão consolidada} \\
\midrule
\endhead
C06 & 2026-07-01 & Exibir dois passos de transformação & Área fixa com ocultação parcial & Aumentar e centralizar a área de diagramas & Renderização adaptada a múltiplos diagramas \\
\addlinespace[2pt]
C10 & 2026-07-02 & Analisar percursos de interação & Redes redundantes e pouco legíveis & Distinguir S/C e organizar redes verticalmente & Redes de transição mais interpretáveis \\
\addlinespace[2pt]
C14 & 2026-07-13 & Reduzir diálogo de entrada numérica & Janela ampla e mensagem extensa & Usar diálogo compacto em duas linhas & Menor carga visual \\
\addlinespace[2pt]
C20 & 2026-07-13 & Evitar explicações excessivas & Campos abertos extensos & Limitar caracteres e habilitar explicação após a escolha & Tela compacta e dados mais comparáveis \\
\addlinespace[2pt]
C22 & 2026-07-13 & Manter alinhamento dos painéis & Coordenadas isoladas & Compartilhar margens e recalcular no redimensionamento & Responsividade consistente \\
\addlinespace[2pt]
C28 & 2026-07-14 & Reforçar a equivalência numérica e melhorar o uso do espaço & Mensagem discreta e distribuição espacial pouco eficiente & Exigir faixa de destaque, cores consistentes para os valores e redistribuição das quatro colunas & Hierarquia visual reforça mesma solução numérica e diferencia valores sem alterar a estrutura conceitual \\
\addlinespace[2pt]
C29 & 2026-07-14 & Explicitar a tese central do Gérard na tela comparativa & Texto secundário com pouco destaque & Substituir a mensagem inferior por uma formulação assertiva e visualmente reforçada & Faixa superior passa a enunciar explicitamente que mesma solução numérica não implica a mesma estrutura matemática \\
\addlinespace[2pt]
C30 & 2026-07-14 & Padronizar componentes visuais auxiliares & Diálogo legado e menus com tipografia e espaçamento inconsistentes & Padronizar diálogo de inserção, fontes, negrito, margens, bordas e realce dos menus & Diálogo e menus passam a compartilhar o padrão visual da interface \\
\addlinespace[2pt]
C31 & 2026-07-14 & Reduzir ruído visual na indicação de submenus & Setas textuais inseridas manualmente ao lado dos itens & Retirar as setas e preservar abertura por onmouseover, realce e clique alternativo & Submenus permanecem descobertos pelo comportamento de hover com menu mais limpo \\
\addlinespace[2pt]
C35 & 2026-07-14 & Representar adequadamente a atividade cognitiva e oferecer apoio sem exposição contínua & Aba denominada Montar e definição permanentemente visível & Substituir montar por construir e mostrar a definição somente no onmouseover do nome da categoria & Nomenclatura reforça raciocínio construtivo e ajuda contextual preserva a leitura limpa do diagrama \\
\addlinespace[2pt]
C38 & 2026-07-15 & Reduzir confusão na leitura do formulário & Expor identificadores técnicos ao pesquisador & Ocultar campos com referência a ID mantendo-os internamente & Curadoria preserva vínculos internos sem expor metadados técnicos \\
\addlinespace[2pt]
C40 & 2026-07-15 & Evitar corte do texto explicativo & Desenhar a descrição em uma linha & Aplicar quebra controlada e posteriormente reduzir informação acessória & Texto deixa de ultrapassar a área disponível \\
\addlinespace[2pt]
C42 & 2026-07-15 & Permitir manipulação direta do valor relativo & Exibir ponto apenas como marcador visual & Tornar o ponto arrastável e restaurar o cartão numérico sincronizado & Eixo, cartão e barras passam a responder à mesma interação \\
\addlinespace[2pt]
C43 & 2026-07-15 & Eliminar sobreposição entre eixo, números e cartão & Posicionar componentes em uma faixa estreita & Ampliar área de captura e separar espacialmente escala e cartão & Representação ganha legibilidade e área de interação mais estável \\
\addlinespace[2pt]
C44 & 2026-07-15 & Manter o eixo estável durante o arraste & Recalcular o máximo com o valor corrente do ponto & Fixar o eixo pelo valor curado e mover o ponto de forma contínua & O eixo não encolhe e o marcador pode percorrer toda a escala \\
\addlinespace[2pt]
C46 & 2026-07-15 & Adequar a área ao conteúdo sinal e número & Usar painel alto semelhante às barras & Reduzir para cartão compacto com valor centralizado & O valor relativo ocupa espaço proporcional à informação exibida \\
\addlinespace[2pt]
C48 & 2026-07-15 & Tornar visíveis os papéis semânticos dos nós & Manter nós sem rótulos contextuais & Adicionar rótulos diretamente aos elementos sem alterar a notação formal & Papéis semânticos passam a ser identificáveis no próprio diagrama \\
\addlinespace[2pt]
C54 & 2026-07-15 & Melhorar uso do espaço vertical & Distribuir unidades com espaçamento amplo & Compactar quadradinhos de baixo para cima e fixar o eixo à geometria das barras & Barras ganham compactação e alinhamento estável \\
\addlinespace[2pt]
C55 & 2026-07-15 & Manter a escala durante a manipulação & Redimensionar a escala conforme o valor corrente & Fixar extremos e espaçamento, movendo somente o ponto de controle & Escala permanece estável durante a exploração \\
\addlinespace[2pt]
C56 & 2026-07-15 & Relacionar papéis e participantes & Exibir apenas os papéis formais & Adicionar subtítulos com personagens ou objetos nos nós & Participantes passam a aparecer junto aos papéis semânticos \\
\addlinespace[2pt]
C57 & 2026-07-15 & Evitar intersecção entre rótulo e seta & Manter o Referendo abaixo do nó superior & Mover o rótulo e o personagem para cima do quadrado & Referendo deixa de se sobrepor à seta vertical \\
\addlinespace[2pt]
C58 & 2026-07-15 & Aprimorar ordem de leitura do nó superior & Exibir Referendo antes do personagem & Colocar o nome do personagem acima do rótulo Referendo & Leitura do participante e do papel fica mais clara \\
\addlinespace[2pt]
C59 & 2026-07-15 & Distinguir personagens dos papéis & Usar a mesma cor para todos os rótulos & Aplicar vermelho suave aos nomes dos personagens em todas as categorias & Personagens ganham identidade cromática consistente \\
\addlinespace[2pt]
C63 & 2026-07-15 & Exibir personagens em todas as categorias & Manter o recurso apenas na comparação & Estender subtítulos e ajustar o espaçamento do Referendo & Personagens passam a ser apresentados em todas as categorias \\
\addlinespace[2pt]
C69 & 2026-07-15 & Melhorar a acomodação da resultante & Usar seta longa e espaço insuficiente para o Todo & Reduzir a seta e redistribuir coleções e valor resultante & Todo permanece legível dentro da área útil \\
\addlinespace[2pt]
C71 & 2026-07-15 & Reduzir informação redundante & Manter texto explicativo abaixo do título & Remover o texto e preservar título e manipulação & Área complementar fica mais limpa \\
\addlinespace[2pt]
C72 & 2026-07-15 & Evitar acesso a funcionalidades incompletas & Manter menus ativos sem implementação consolidada & Desabilitar opções e informar Em construção & Usuário deixa de entrar em fluxos ainda não disponíveis \\
\addlinespace[2pt]
C74 & 2026-07-15 & Facilitar o envio do registro & Salvar somente a cópia local & Abrir cliente de e-mail com mensagem preenchida & Relato pode ser enviado sem transcrição manual \\
\addlinespace[2pt]
C76 & 2026-07-15 & Atualizar fotografia do submenu & Manter imagem anterior & Substituir a foto preservando redimensionamento automático & Submenu passa a exibir a nova imagem \\
\addlinespace[2pt]
C77 & 2026-07-15 & Exibir duas fotografias & Mostrar somente uma imagem & Organizar as fotos lado a lado com escala automática & Duas imagens passam a coexistir no submenu \\
\addlinespace[2pt]
C78 & 2026-07-15 & Substituir uma fotografia preservando o layout & Alterar a imagem com risco de deformação & Manter proporção e alinhamento com a segunda foto & Nova foto mantém o equilíbrio do painel \\
\addlinespace[2pt]
C79 & 2026-07-15 & Contextualizar o registro histórico & Usar o rótulo Em Recife e legenda inferior & Renomear para Na UFPE em 2009 e retirar a legenda & Submenu passa a indicar local e data de forma direta \\
\addlinespace[2pt]
C80 & 2026-07-15 & Atualizar destinatário & Enviar para endereço institucional anterior & Usar anaemilia@gmail.com mantendo cópia local & Relatos passam a ser dirigidos ao endereço definido \\
\addlinespace[2pt]
C81 & 2026-07-15 & Sinalizar que o item abre outro painel & Usar apenas o texto & Adicionar seta discreta e ajustar o posicionamento & Item passa a indicar continuidade hierárquica \\
\addlinespace[2pt]
C82 & 2026-07-15 & Adequar a posição da seta & Colocar o indicador antes do texto & Mover a seta para depois do rótulo & Indicador segue o padrão de submenu \\
\addlinespace[2pt]
C83 & 2026-07-15 & Evitar glifo ausente & Usar caractere Unicode dependente da fonte & Desenhar a seta como ícone vetorial Swing & Seta deixa de aparecer como quadrado substituto \\
\addlinespace[2pt]
C84 & 2026-07-15 & Separar texto e indicador & Posicionar a seta junto ao rótulo & Usar toda a largura com texto à esquerda e seta à direita & Item ganha alinhamento estável e legível \\
\addlinespace[2pt]
C89 & 2026-07-15 & Preencher efetivamente a mensagem no navegador & Usar mailto associado ao Gmail & Abrir o compositor Web com destinatário, assunto e corpo parametrizados & Gmail passa a abrir com o relato preenchido \\
\addlinespace[2pt]
C97 & 2026-07-16 & Integrar a janela à identidade visual do Gérard & Manter o formulário com o tema cinza padrão do Swing & Aplicar fundo claro, cartões brancos, títulos azuis e botões consistentes sem alterar os dados & A janela qualitativa passa a pertencer visualmente à mesma atividade educacional \\
\addlinespace[2pt]
C103 & 2026-07-16 & Evitar sobreposição entre controles e números de cardinalidade & Posicionar os ícones na mesma faixa lateral dos números & Separar espacialmente controles e contagens sem afastá-los da representação & Controles permanecem associados ao agrupamento sem cobrir os numerais \\
\addlinespace[2pt]
C104 & 2026-07-16 & Aumentar a folga visual entre controles e contagens & Considerar suficiente a primeira correção de sobreposição & Refinar novamente os afastamentos horizontal e vertical & Coluna de controles passa a ocupar a borda superior lateral com separação confortável \\
\addlinespace[2pt]
C109 & 2026-07-16 & Uniformizar a ênfase dos personagens & Manter pesos tipográficos diferentes entre representações & Exibir em negrito todos os nomes curados sem alterar cor, posição ou papel & Personagens recebem a mesma hierarquia tipográfica em todos os diagramas \\
\addlinespace[2pt]
C110 & 2026-07-16 & Comunicar que os quadradinhos são manipuláveis & Manter unidades totalmente planas ou aplicar volume forte & Usar relevo 2,5D discreto, intensificado apenas no foco e no arraste & Unidades permanecem contáveis e ganham affordance sem alterar a paleta semântica \\
\addlinespace[2pt]
C113 & 2026-07-16 & Tornar perceptível o momento em que um objeto é pego & Alterar apenas o cursor ou manter a pintura original & Combinar cursor, escala de 106\%, elevação, sombra e prioridade de pintura & Elementos em trânsito recebem prioridade visual sem alterar coordenadas reais \\
\addlinespace[2pt]
C114 & 2026-07-16 & Padronizar a mão usada nos elementos manipuláveis & Usar cursores vetoriais personalizados & Adotar o cursor nativo já consolidado nos controles laterais & A mesma mão passa a indicar disponibilidade em toda a interface \\
\addlinespace[2pt]
C115 & 2026-07-16 & Distinguir disponibilidade de movimento ativo & Manter a mesma mão durante todas as fases & Usar mão no mouseover e cursor de movimento durante o arraste & Cursor comunica separadamente que o item pode ser pego e que está em deslocamento \\
\addlinespace[2pt]
C121 & 2026-07-16 & Explicitar massa e inércia no seguimento elástico & Usar posição suavizada sem modelo físico explícito & Aplicar massa-mola-amortecedor com limite de velocidade, atraso máximo e soltura exata & O elemento persegue o cursor com pequeno atraso e momento sem afetar regras da C120 \\
\addlinespace[2pt]
\bottomrule
\end{longtable}
\endgroup
\end{landscape}